\documentclass[12pt,a4paper]{article}

%==================== 宏包 ====================
\usepackage[UTF8,fontset=windows]{ctex}
\usepackage{amsmath,amssymb,amsthm}
\usepackage{mathtools}
\usepackage{geometry}
\usepackage{enumitem}
\usepackage{xcolor}
\usepackage{tcolorbox}
\usepackage{booktabs}
\usepackage{array}
\usepackage{hyperref}
\usepackage{fancyhdr}
\usepackage{titlesec}
\usepackage{algorithm}
\usepackage{algpseudocode}
\usepackage{setspace}

%==================== 页面设置 ====================
\geometry{left=2.5cm,right=2.5cm,top=2.5cm,bottom=2.5cm}

%==================== 字体设置 ====================
\setCJKmainfont{SimSun}[AutoFakeBold=2]  % 宋体
\setCJKsansfont{SimHei}[AutoFakeBold=2]  % 黑体
\setmainfont{Times New Roman}

%==================== 标题格式设置 ====================
% 章标题：黑体16磅加粗居中，段前24磅，段后18磅
\titleformat{\section}
{\centering\fontsize{16pt}{20pt}\selectfont\CJKfamily{hei}\bfseries}
{\thesection\quad}
{0em}
{}
\titlespacing*{\section}{0pt}{24pt}{18pt}

% 一级标题：黑体14磅加粗居中，段前24磅，段后6磅
\titleformat{\subsection}
{\centering\fontsize{14pt}{18pt}\selectfont\CJKfamily{hei}\bfseries}
{\thesubsection\quad}
{0em}
{}
\titlespacing*{\subsection}{0pt}{24pt}{6pt}

% 二级标题：黑体13磅居左，段前12磅，段后6磅
\titleformat{\subsubsection}
{\fontsize{13pt}{16pt}\selectfont\CJKfamily{hei}\bfseries}
{\thesubsubsection\quad}
{0em}
{}
\titlespacing*{\subsubsection}{0pt}{12pt}{6pt}

%==================== 段落格式设置 ====================
\setlength{\parindent}{2em}  % 首行缩进2汉字符
\setlength{\parskip}{0pt}    % 段间距0磅
\linespread{1.25}            % 固定值行距20磅（约1.25倍）

%==================== 自定义环境 ====================
\newtcolorbox{keybox}[1][]{
  colback=yellow!8!white,
  colframe=orange!60!black,
  fonttitle=\bfseries,
  title=关键点,
  #1
}
\newtcolorbox{notebox}[1][]{
  colback=blue!5!white,
  colframe=blue!55!black,
  fonttitle=\bfseries,
  title=备注,
  #1
}
\newtcolorbox{summarybox}[1][]{
  colback=green!6!white,
  colframe=green!50!black,
  fonttitle=\bfseries,
  title=小结,
  #1
}

%==================== 笔记空间 ====================
\newcommand{\notespace}[1][3cm]{%
  \vspace{#1}%
  \noindent\textbf{【笔记空间】}%
  \vspace{#1}%
}

%==================== 数学命令 ====================
\newcommand{\E}{\mathbb{E}}
\newcommand{\R}{\mathbb{R}}
\newcommand{\1}{\mathbf{1}}
\newcommand{\argmax}{\operatorname*{arg\,max}}
\newcommand{\argmin}{\operatorname*{arg\,min}}

%==================== 页眉页脚 ====================
\pagestyle{fancy}
\fancyhf{}
\fancyhead[L]{\textbf{强化学习复习笔记}}
\fancyhead[R]{\thepage}
\renewcommand{\headrulewidth}{0.4pt}

%==================== 文档信息 ====================
\title{\textbf{强化学习复习笔记}\\[0.5em]\large 基于课程内容整理}
\author{}
\date{}

\begin{document}
\maketitle
\tableofcontents
\newpage

%==================================================================
% 第一章 强化学习概述
%==================================================================
\section{强化学习概述}

\subsection{强化学习能解决的问题}

强化学习适合解决序贯决策、复杂控制和交互优化问题，典型应用包括：

\begin{itemize}[leftmargin=2em]
  \item 非线性控制、视频游戏、机器人控制（走路、站立）
  \item AlphaGo、多智能体协作（星际争霸）
  \item 海量智能体问题、可控核聚变控制
  \item 大语言模型中的人类反馈优化（如 ChatGPT 的 RLHF）
\end{itemize}

\subsubsection*{典型场景：可控核聚变}

托卡马克装置利用电磁力约束核聚变时的超高温等离子体。关键点：

\begin{itemize}[leftmargin=2em]
  \item 控制等离子体形状，服务于发电
  \item 热传导不能过多（容器融化）也不能过少（发电效率低）
\end{itemize}

\subsubsection*{典型场景：ChatGPT 训练思路}

\begin{enumerate}[leftmargin=2em]
  \item 收集示例数据，得到有监督微调策略 SFT
  \item 收集人工标注的对比数据，训练回报模型
  \item 使用 PPO 算法和学到的回报模型进一步优化 SFT
\end{enumerate}

\notespace[2cm]

\subsection{强化学习与监督学习的联系和区别}

\textbf{相同点：} 都从数据中学习；都会逐步改善性能。

\textbf{不同点：}

\begin{itemize}[leftmargin=2em]
  \item \textbf{数据类型}：监督学习需要标签数据；强化学习需要与环境交互得到的数据
  \item \textbf{优化目标}：监督学习优化预测误差；强化学习优化长期累积回报
\end{itemize}

\notespace[2cm]

\subsection{强化学习与其他优化方法}

序贯最优决策的其他方法：

\begin{itemize}[leftmargin=2em]
  \item 状态空间很小时：线性规划（资源分配、调度）
  \item 模型已知、回报函数解析时：动态规划或最优控制
  \item 控制策略较简单时：进化算法（遗传算法）
\end{itemize}

\subsubsection*{最优控制的一般形式}

状态方程：
\begin{equation}
\dot{X} = f(t, X, U), \quad X(t_0)=X_0
\end{equation}

性能指标函数：
\begin{equation}
J(x,t)=\phi(X(t_f),t_f)+\int_t^{t_f} L(x(\tau),u(\tau),\tau)\,d\tau
\end{equation}

最优值函数：
\begin{equation}
V(X,t)=\min_{u\in\Omega}\left[\phi(X(t_f),t_f)+\int_t^{t_f} L(x(\tau),u(\tau),\tau)\,d\tau\right]
\end{equation}

对应 HJB 方程：
\begin{equation}
-\frac{\partial V}{\partial t} = \min_{u\in U} \left\{ L(x,t,u) + \left(\frac{\partial V}{\partial x}\right)^{\!T} f(x,t,u) \right\}
\end{equation}

\begin{keybox}
当模型未知或非常复杂、回报函数未知时，强化学习通过采样和环境交互来学习最优策略。
\end{keybox}

\notespace[2cm]

\subsection{基本概念}

\subsubsection*{策略（Policy）}

策略是从状态空间到动作空间的映射：
\begin{equation}
\pi(\cdot \mid s): s \rightarrow a
\end{equation}

\subsubsection*{奖赏（Reward）}

环境在每个时间步给智能体的立即回报。

\subsubsection*{值函数（Value Function）}

从某个状态开始，按照某个策略行动时所能获得的累积回报。

\subsubsection*{学习目标}

使长期累积回报最大，即让值函数最大。

\notespace[2cm]

\subsection{与环境交互的学习过程}

强化学习通过与环境交互产生数据流，智能体利用这些数据学习并优化行为，通常放在马尔可夫决策过程（MDP）框架下理解。

强化学习的两个核心特征：

\begin{itemize}[leftmargin=2em]
  \item \textbf{试错探索}
  \item \textbf{延迟回报}
\end{itemize}

\subsubsection*{井字游戏示例}

\begin{itemize}[leftmargin=2em]
  \item 设置值函数列表，每个值表示从当前状态出发最终获胜的概率
  \item 初始化：赢的状态取 $1$，输的状态取 $0$，其他状态取 $0.5$
  \item 价值更新（备份更新）：
\end{itemize}

\begin{equation}
V(S_t) \leftarrow V(S_t) + \alpha \bigl[V(S_{t+1}) - V(S_t)\bigr]
\end{equation}

\notespace[2cm]

\subsection{强化学习算法历史与发展}

\subsubsection*{发展阶段}

\begin{itemize}[leftmargin=2em]
  \item 1998 年前：基本理论框架形成
  \item 1998--2013 年：各种直接策略搜索方法出现
  \item 2013 年至今：深度强化学习快速发展
\end{itemize}

\subsubsection*{代表方向}

\begin{itemize}[leftmargin=2em]
  \item 基于值函数的方法：DQN、DNC、AlphaGo
  \item 基于策略梯度的强化学习
  \item 基于 EM、基于路径积分、基于回归、基于模型的强化学习
\end{itemize}

\notespace[2cm]

\subsection{强化学习的分类}

\subsubsection*{按是否依赖模型}

\begin{itemize}[leftmargin=2em]
  \item \textbf{基于模型}：先学习系统模型，再基于模型求最优策略
  \item \textbf{无模型}：不显式学习环境模型，直接通过交互数据得到最优策略
\end{itemize}

\subsubsection*{按策略更新方式}

\begin{itemize}[leftmargin=2em]
  \item \textbf{基于值函数}：先求最优值函数，再由值函数重构最优策略
  \item \textbf{基于直接策略搜索}：直接在策略空间中搜索最优策略
  \item \textbf{Actor-Critic}：同时逼近值函数和最优策略，类似策略迭代
\end{itemize}

\subsubsection*{按回报函数是否已知}

\begin{itemize}[leftmargin=2em]
  \item \textbf{正向强化学习}：从回报中学习最优策略
  \item \textbf{逆向强化学习}：从专家示例中学习回报函数
\end{itemize}

\subsubsection*{按任务形式和规模}

分层强化学习、元强化学习、多智能体强化学习、迁移学习。

\notespace[2cm]

\subsection{强化学习的特点}

\begin{keybox}
核心在于：搜索/探索 + 正确的评估。利用对当前动作的评估进行学习。
\end{keybox}

\notespace[3cm]

\newpage

%==================================================================
% 第二章 数学基础
%==================================================================
\section{数学基础}

\subsection{概率论基础}

\subsubsection*{期望与方差}

离散型随机变量期望：
\begin{equation}
\E[X]=\sum_x x\,P(X=x)
\end{equation}

连续型随机变量期望：
\begin{equation}
\E[X]=\int x f(x)\,dx
\end{equation}

方差：
\begin{equation}
\mathrm{Var}(X)=\E\bigl[(X-\E[X])^2\bigr]
\end{equation}

\subsubsection*{条件概率}

在事件 $A$ 已发生条件下，事件 $B$ 发生的概率：$P(B\mid A)$。

\subsubsection*{随机过程}

设 $(\Omega,\Sigma,P)$ 是概率空间，对每个参数 $t\in T$，$X(t,\omega)$ 是定义在该概率空间上的随机变量，则随机变量族
\begin{equation}
\{X_t=X(t,\omega);\, t\in T\}
\end{equation}
称为该概率空间上的一个随机过程。

\notespace[2cm]

\subsection{马尔可夫性}

\begin{keybox}
\textbf{马尔可夫性}：系统的下一个状态只与当前状态有关，而与更早的历史无关。
\begin{equation}
P(S_{t+1}\mid S_t)=P(S_{t+1}\mid S_1,\dots,S_t)
\end{equation}
\end{keybox}

\textbf{直观理解：} 当前状态蕴含了所有相关的历史信息；一旦当前状态已知，历史信息就可以被抛弃。

\notespace[3cm]

\newpage

%==================================================================
% 第三章 多臂赌博机问题
%==================================================================
\section{多臂赌博机问题}

\subsection{基本概念}

\textbf{多臂赌博机（Multi-Armed Bandit）}：

\begin{itemize}[leftmargin=2em]
  \item 有 $K$ 个臂
  \item 每次摇动某个臂，会得到服从该臂对应概率分布的回报
  \item 目标：通过学习知道摇哪个臂能获得更高回报
\end{itemize}

\textbf{问题目标：} 在摇动 $N$ 次的过程中，使总回报最高。

\subsubsection*{强化学习视角下的表述}

\begin{itemize}[leftmargin=2em]
  \item 动作：$a \in \{0,1,2,\dots\}$
  \item 状态：单一状态 $s$
  \item 回报：$r$，服从不同动作对应的不同回报分布
  \item 学习目标：在状态 $s$ 下找到最优动作，利用回报 $r$ 学习动作的优劣
\end{itemize}

\notespace[2cm]

\subsection{动作价值的定义}

动作的真实价值：
\begin{equation}
q_*(a)=\E[R_t \mid A_t=a]
\end{equation}

动作的估计值（样本平均）：
\begin{equation}
Q_t(a)=\frac{\sum_{i=1}^{t-1} R_i \cdot \1(A_i=a)}{\sum_{i=1}^{t-1} \1(A_i=a)}
\end{equation}
其中 $\1(A_i=a)$ 表示当第 $i$ 次选择动作 $a$ 时取 1，否则取 0。

\notespace[2cm]

\subsection{探索与利用}

在每个臂都至少尝试一次后，选择下一个臂需要平衡：

\begin{itemize}[leftmargin=2em]
  \item \textbf{利用（exploitation）}：选择估计值最大的动作
\end{itemize}

\begin{equation}
a=\argmax_a q(a)
\end{equation}

\begin{itemize}[leftmargin=2em]
  \item \textbf{探索（exploration）}：选择非贪婪动作，以获取更多信息
\end{itemize}

\notespace[2cm]

\subsection{增量式更新}

\begin{keybox}
\textbf{新估计值 = 旧估计值 + 步长 $\times$ 预测误差}
\begin{equation}
Q_{n+1} = Q_n + \frac{1}{n}\bigl(R_n - Q_n\bigr)
\end{equation}
\end{keybox}

\notespace[2cm]

\subsection{常见动作选择方法}

\subsubsection*{Upper Confidence Bound（UCB）}

\begin{equation}
A_t=\argmax_a \left[ Q_t(a)+c\sqrt{\frac{\ln t}{N_t(a)}} \right]
\end{equation}
其中 $Q_t(a)$ 为价值估计，$N_t(a)$ 为动作 $a$ 被选择的次数，$c$ 控制探索强度。

\subsubsection*{玻尔兹曼分布（Softmax）动作选择}

\begin{equation}
\Pr\{A_t=a\}=\frac{e^{H_t(a)}}{\sum_{b=1}^{k} e^{H_t(b)}}
\end{equation}
其中 $H_t(a)$ 表示动作偏好（preference）。

\notespace[2cm]

\subsection{从多臂赌博机到马尔可夫决策过程}

\begin{notebox}
\textbf{核心区别：}
\begin{itemize}[leftmargin=1.5em]
  \item 多臂赌博机：单状态问题，动作只影响立即回报
  \item MDP：多状态问题，动作不仅影响回报，还影响状态转移
\end{itemize}
\end{notebox}

MDP 中：在状态 $s_t$ 执行动作 $a_t$，获得回报 $r_t$，转移到下一状态 $s_{t+1}$：
\begin{equation}
s_t \xrightarrow{a_t,\, r_t} s_{t+1}
\end{equation}

\textbf{延迟回报：} 在 MDP 中，动作的好坏不能只看立即回报，还要看未来回报。需要平衡立即回报与将来回报。

\notespace[3cm]

\newpage

%==================================================================
% 第四章 马尔可夫决策过程
%==================================================================
\section{马尔可夫决策过程}

\subsection{定义}

\begin{keybox}
MDP 由五元组描述：$(S,A,P,R,\gamma)$
\begin{itemize}[leftmargin=1.5em]
  \item $S$：有限状态集
  \item $A$：有限动作集
  \item $P$：状态转移概率
  \item $R$：回报函数
  \item $\gamma$：折扣因子
\end{itemize}
\end{keybox}

\notespace[2cm]

\subsection{动力学模型}

随机变量 $s',r$ 的联合分布：
\begin{equation}
p(s',r\mid s,a)=\Pr\{S_t=s',R_t=r\mid S_{t-1}=s,A_{t-1}=a\}
\end{equation}

\subsubsection*{状态转移概率}

\begin{equation}
p(s'\mid s,a)=\sum_r p(s',r\mid s,a)
\end{equation}

\subsubsection*{回报的期望}

\begin{equation}
R(s,a)=\E[R_t\mid S_{t-1}=s,A_{t-1}=a] =\sum_r \sum_{s'} r\, p(s',r\mid s,a)
\end{equation}

\notespace[2cm]

\subsection{状态、动作与回报}

\subsubsection*{状态}

状态是任何我们能够知道、并且可能有助于决策的信息。

\begin{itemize}[leftmargin=2em]
  \item DQN 的状态：连续的四帧图像
  \item AlphaGo 的状态：过去八手的双方棋局
\end{itemize}

\textbf{状态表示会严重影响学习效果，目前更像一门艺术。}

\subsubsection*{动作}

动作是任何我们希望智能体学习做出的决策，可以是宏观/微观、精神/物质层面动作。

\subsubsection*{立即回报的设计}

立即回报的设置应服务于最终目标：

\begin{itemize}[leftmargin=2em]
  \item 机器人学习走路：回报正比于前进速度
  \item 机器人逃离迷宫：每走一步给 $-1$，鼓励尽快逃离
  \item 下棋：赢 $+1$，输 $-1$，平局 $0$
\end{itemize}

\notespace[2cm]

\subsection{累积回报与折扣回报}

\subsubsection*{累积回报}

\begin{equation}
G_t = R_{t+1}+R_{t+2}+\cdots+R_T
\end{equation}

\subsubsection*{幕与任务类型}

\begin{itemize}[leftmargin=2em]
  \item \textbf{幕（episode）}：一次完整交互序列
  \item \textbf{终止状态（terminal state）}：交互结束的状态
  \item \textbf{episodic tasks}：有终止状态的任务
  \item \textbf{continuing tasks}：没有终止状态、持续进行的任务
\end{itemize}

\subsubsection*{折扣回报}

\begin{equation}
G_t = R_{t+1}+\gamma R_{t+2}+\gamma^2 R_{t+3}+\cdots = \sum_{k=0}^{\infty}\gamma^k R_{t+k+1}
\end{equation}
其中 $0 \le \gamma \le 1$。

\begin{keybox}
\textbf{折扣率 $\gamma$ 的含义：}
\begin{itemize}[leftmargin=1.5em]
  \item $\gamma \to 1$：智能体更有远见（farsighted），重视未来回报
  \item $\gamma=0$：智能体短视（myopic），只最大化立即回报
  \item 若回报有界：$G_t \le \dfrac{R_{\max}}{1-\gamma}$
\end{itemize}
\end{keybox}

\notespace[2cm]

\subsection{策略与值函数}

\subsubsection*{策略的定义}

策略 $\pi$ 是给定状态 $s$ 时在动作集上的一个概率分布：
\begin{equation}
\pi(a\mid s)=P(A_t=a\mid S_t=s)
\end{equation}

\subsubsection*{常见策略}

\begin{enumerate}[leftmargin=2em]
  \item \textbf{贪婪策略}：
\end{enumerate}

\begin{equation}
\pi(a\mid s)= \begin{cases} 1, & a=\argmax_{a\in A} q_*(s,a) \\ 0, & \text{otherwise} \end{cases}
\end{equation}

\begin{enumerate}[leftmargin=2em]
  \setcounter{enumi}{1}
  \item \textbf{随机策略}：按某个概率分布随机选动作
  \item \textbf{高斯策略}：常用于连续动作空间
  \item \textbf{玻尔兹曼策略}：
\end{enumerate}

\begin{equation}
\pi_\theta(a\mid s)= \frac{\exp(Q(s,a;\theta))} {\sum_b \exp(Q(s,b;\theta))}
\end{equation}

\notespace[2cm]

\subsection{值函数与动作价值函数}

\subsubsection*{状态值函数}

\begin{equation}
v_\pi(s)=\E_\pi[G_t\mid S_t=s] =\E_\pi\left[\sum_{k=0}^{\infty}\gamma^k R_{t+k+1}\,\middle|\, S_t=s\right]
\end{equation}

\subsubsection*{行为值函数}

\begin{equation}
q_\pi(s,a)=\E_\pi[G_t\mid S_t=s,A_t=a] =\E_\pi\left[\sum_{k=0}^{\infty}\gamma^k R_{t+k+1}\,\middle|\, S_t=s,A_t=a\right]
\end{equation}

\begin{notebox}
由于状态转移概率和策略都可能是随机的，因此折扣累积回报 $G_t$ 本身是随机变量。值函数就是该随机回报的期望。
\end{notebox}

\notespace[2cm]

\subsection{贝尔曼方程}

\subsubsection*{值函数与动作价值函数之间的关系}

\begin{equation}
v_\pi(s)=\sum_{a\in A}\pi(a\mid s)\,q_\pi(s,a)
\end{equation}

\subsubsection*{状态值函数的贝尔曼方程}

\begin{equation}
v_\pi(s)=\sum_a \pi(a\mid s)\sum_{s',r} p(s',r\mid s,a)\bigl[r+\gamma v_\pi(s')\bigr]
\end{equation}

\subsubsection*{行为值函数的贝尔曼方程}

\begin{equation}
q_\pi(s,a)=\sum_{s',r} p(s',r\mid s,a)\bigl[r+\gamma \sum_{a'}\pi(a'\mid s')q_\pi(s',a')\bigr]
\end{equation}

\subsubsection*{最优贝尔曼方程}

\begin{equation}
v^*(s)=\max_a \sum_{s',r} p(s',r\mid s,a)\bigl[r+\gamma v^*(s')\bigr]
\end{equation}

\begin{equation}
q^*(s,a)=\sum_{s',r} p(s',r\mid s,a)\bigl[r+\gamma \max_{a'} q^*(s',a')\bigr]
\end{equation}

\notespace[3cm]

\newpage

%==================================================================
% 第五章 动态规划
%==================================================================
\section{动态规划}

\subsection{动态规划的本质}

\begin{keybox}
动态规划的核心思想：将多阶段决策问题通过贝尔曼方程，转化为多个单阶段决策问题。
\end{keybox}

\subsubsection*{离散情形}

\begin{equation}
J^*(x(j),j) = \min_{u(j)\in U} \left\{ L(x(j),u(j),j)+J^*(x(j+1),j+1) \right\}
\end{equation}
其中 $x(j+1)=f(x(j),u(j),j)$。

\subsubsection*{连续情形与 HJB 方程}

\begin{equation}
-\frac{\partial J^*[x(t),t]}{\partial t} = \min_{u(t)\in U} \left\{ L[x(t),u(t),t] + \frac{\partial J^*[x(t),t]}{\partial x^T}f[x(t),u(t),t] \right\}
\end{equation}

\notespace[2cm]

\subsection{策略评估}

\textbf{输入：} 策略 $\pi$、状态转移概率 $P_{ss'}^a$、回报函数 $R_{sa}$、折扣因子 $\gamma$

\textbf{初始化：} $V(s)=0$

\textbf{迭代更新：}

\begin{equation}
v_{k+1}(s) = \sum_{a\in A}\pi(a|s)\left(R_{sa} + \gamma\sum_{s'\in S}P_{ss'}^a v_k(s')\right)
\end{equation}

直到 $v_{k+1} = v_k$。该方法属于"期望更新"，一次对所有状态的更新称为"一次状态扫描"。

\notespace[2cm]

\subsection{策略改进}

\textbf{基本思想：} 在每个状态上采用贪婪策略：

\begin{equation}
\pi_{l+1}(s)\in \argmax_a q_{\pi_l}(s,a)
\end{equation}

\textbf{策略改善理论：} 如果新策略在每个状态上的动作选择都不差于原策略，则新策略不会比原策略差，通常会更好。

\textbf{贪婪策略计算：}

\begin{equation}
q_{\pi_l}(s,a)=R_{sa}+\gamma\sum_{s'\in S}P_{ss'}^a v_{\pi_l}(s')
\end{equation}

\begin{equation}
\pi_{l+1}(s)\in \argmax_a q_{\pi_l}(s,a)
\end{equation}

\notespace[2cm]

\subsection{策略迭代}

\begin{keybox}
策略迭代由两步交替进行，直到策略不再变化：
\begin{enumerate}[leftmargin=1.5em]
  \item \textbf{策略评估}：对当前策略 $\pi_l$ 求得 $V^{\pi_l}$
  \item \textbf{策略改进}：$\pi_{l+1}(s)\in \argmax_a q_{\pi_l}(s,a)$
\end{enumerate}
\end{keybox}

\textbf{终止条件：} $\pi_{l+1}=\pi_l$，输出 $\pi^*=\pi_l$。

\textbf{关系总结：} $V^* \ge V^{\pi}$，当 $\pi = \text{greedy}(V)$ 时，通过反复评估和改进可逐步逼近最优值函数 $V^*$ 和最优策略 $\pi^*$。

\notespace[2cm]

\subsection{值函数迭代}

\textbf{动机：} 策略改进一定要等到值函数完全收敛吗？不一定。当 $K=1$ 时便进行策略改进，可得到值迭代算法。

\textbf{最优贝尔曼方程：}

\begin{equation}
v^*(s)=\max_a\left( R_{sa}+\gamma\sum_{s'\in S}P_{ss'}^a v^*(s') \right)
\end{equation}

\textbf{算法：}

\begin{algorithmic}[1]
\State 初始化 $v(s)=0$
\Repeat
  \For{每个 $s$}
    \State $v_{l+1}(s) = \max_a \left(R_{sa} + \gamma \sum_{s'} P_{ss'}^a v_l(s')\right)$
  \EndFor
\Until{$v_{l+1} = v_l$}
\end{algorithmic}

\textbf{收敛后提取策略：}

\begin{equation}
\pi(s)=\argmax_a\left( R_{sa}+\gamma\sum_{s'\in S}P_{ss'}^a v_l(s') \right)
\end{equation}

\begin{summarybox}
\textbf{特点：} 将策略评估与策略改进更紧密地结合，不需要每次都把策略评估做到完全收敛，通常比完整策略迭代更直接。
\end{summarybox}

\notespace[3cm]

\newpage

%==================================================================
% 第六章 蒙特卡洛方法
%==================================================================
\section{蒙特卡洛方法}

\subsection{基本思想}

\begin{keybox}
蒙特卡洛方法的核心思想：利用经验平均代替随机变量的期望。
\end{keybox}

一次实验（episode）是从初始状态开始，到终止状态结束的一条完整轨迹。对某个时刻 $t$ 之后的折扣回报：

\begin{equation}
G_t=R_{t+1}+\gamma R_{t+2}+\gamma^2 R_{t+3}+\cdots+\gamma^{T-t-1}R_T
\end{equation}

\notespace[2cm]

\subsection{与动态规划策略评估对比}

\begin{itemize}[leftmargin=2em]
  \item \textbf{DP}：需要已知模型，每次进行一次状态扫描
  \item \textbf{MC}：每次进行一次完整试验，用试验中观察到的回报估计值函数
\end{itemize}

\subsection{First-visit MC 与 Every-visit MC}

\begin{itemize}[leftmargin=2em]
  \item \textbf{First-visit MC}：每幕中只使用某状态第一次出现后的回报进行估计
  \item \textbf{Every-visit MC}：每幕中该状态每次出现后的回报都用于估计
\end{itemize}

根据大数定律：$\hat{v}(s)\to v_{\pi}(s)$，当 $N(s)\to\infty$。

\notespace[2cm]

\subsection{探索型初始化}

为保证充分访问，采用探索型初始化（exploring starts）：每个状态都有一定几率作为初始状态；更强的形式是每个状态-行为对都有机会作为初始状态-动作对出现。

\notespace[2cm]

\subsection{蒙特卡洛评估行为值函数}

在模型未知时，最重要的是评估行为值函数：

\begin{equation}
q_{\pi}(s,a)=\E_{\pi}\left[\sum_{k=0}^{\infty}\gamma^k R_{t+k+1}\mid S_t=s,A_t=a\right]
\end{equation}

\begin{notebox}
\textbf{蒙特卡洛评估与动态规划评估的区别：}
\begin{itemize}[leftmargin=1.5em]
  \item MC 不需要 bootstrap
  \item MC 对单个状态的估计相对独立
  \item MC 对单个状态的估计复杂度不直接依赖于整个状态空间大小
  \item DP 依赖模型；MC 依赖采样得到的完整 episode
\end{itemize}
\end{notebox}

\notespace[2cm]

\subsection{蒙特卡洛策略改进}

贪婪改进形式：

\begin{equation}
\pi(s)=\argmax_a q_{\pi}(s,a)
\end{equation}

\subsubsection*{蒙特卡洛强化学习：探索初始化}

\begin{algorithmic}[1]
\State 初始化 $Q(s,a)$、$\pi(s)$、$\mathrm{Returns}(s,a)$
\Repeat
  \State 随机选择初始状态-动作对 $(S_0,A_0)$
  \State 从 $(S_0,A_0)$ 出发，按策略 $\pi$ 生成一个 episode
  \For{episode 中每个 $(s,a)$}
    \State 计算 $G$（第一次出现后的回报）
    \State 将 $G$ 附加到 $\mathrm{Returns}(s,a)$
    \State $Q(s,a)\leftarrow \mathrm{average}(\mathrm{Returns}(s,a))$
  \EndFor
  \State $\pi(s)\leftarrow \argmax_a Q(s,a)$
\Until{收敛}
\end{algorithmic}

\notespace[2cm]

\subsection{On-policy 与 Off-policy}

\subsubsection*{基本要求}

策略必须是永久温和的，即对任意状态 $s$ 和动作 $a$：$\pi(a|s)>0$。

\subsubsection*{On-policy MC}

产生数据的策略、被评估的策略、被改进的策略三者是同一个策略。典型做法是使用 $\varepsilon$-soft 策略：

\begin{equation}
\pi(a|s)= \begin{cases} 1-\varepsilon+\dfrac{\varepsilon}{|A(s)|}, & a=a^* \\ \dfrac{\varepsilon}{|A(s)|}, & a\neq a^* \end{cases}
\end{equation}

\subsubsection*{Off-policy MC}

使用两个不同的策略：

\begin{itemize}[leftmargin=2em]
  \item \textbf{目标策略} $\pi$：要学习和改进的策略，通常较贪婪
  \item \textbf{行为策略} $\mu$：用来实际采样，负责探索
\end{itemize}

\textbf{覆盖性条件：} 若 $\pi(a|s)>0$，则必须有 $\mu(a|s)>0$。

\notespace[2cm]

\subsection{重要性采样}

\subsubsection*{普通重要性采样}

\begin{equation}
\E[f]\approx \frac{1}{N}\sum_{n=1}^{N}\omega_n f(z_n), \quad \omega_n=\frac{p(z_n)}{q(z_n)}
\end{equation}
特点：无偏估计，但方差可能非常大。

\subsubsection*{加权重要性采样}

\begin{equation}
\E[f]\approx \frac{\sum_{n=1}^{N}\omega_n f(z_n)}{\sum_{n=1}^{N}\omega_n}
\end{equation}
特点：一般有偏，但方差更小、更稳定。

\subsubsection*{MC 中的重要性采样}

重要性比率只与动作概率有关：

\begin{equation}
\rho_{t:T-1} =\prod_{k=t}^{T-1}\frac{\pi(A_k|S_k)}{\mu(A_k|S_k)}
\end{equation}

\textbf{普通重要性采样估计：}

\begin{equation}
V(s)=\frac{\sum_{t\in T(s)} \rho_t G_t}{|T(s)|}
\end{equation}

\textbf{加权重要性采样估计：}

\begin{equation}
V(s)=\frac{\sum_{t\in T(s)} \rho_t G_t}{\sum_{t\in T(s)} \rho_t}
\end{equation}

\notespace[3cm]

\newpage

%==================================================================
% 第七章 时间差分学习
%==================================================================
\section{时间差分学习}

\subsection{TD 学习的基本思想}

\begin{keybox}
\begin{itemize}[leftmargin=1.5em]
  \item \textbf{TD 学习}是采样更新
  \item \textbf{DP 学习}是期望更新
\end{itemize}
\end{keybox}

\subsubsection*{DP、MC、TD 的对比}

\begin{itemize}[leftmargin=2em]
  \item \textbf{MC}：利用完整回报 $G_t$ 估计值函数（采样平均逼近期望）
  \item \textbf{TD}：结合 MC 和 DP，采样得到目标，并利用当前估计值 $V(S_{t+1})$ 进行自举
  \item \textbf{DP}：依赖环境模型提供期望，也利用当前估计值进行更新
\end{itemize}

\notespace[2cm]

\subsection{增量式 MC 与 TD(0)}

\textbf{增量式 MC：}

\begin{equation}
V(S_t)\leftarrow V(S_t)+\alpha\bigl[G_t-V(S_t)\bigr]
\end{equation}

\textbf{TD(0)：}

\begin{equation}
V(S_t)\leftarrow V(S_t)+\alpha\bigl[R_{t+1}+\gamma V(S_{t+1})-V(S_t)\bigr]
\end{equation}
其中 $R_{t+1}+\gamma V(S_{t+1})$ 称为 TD 目标。

\begin{keybox}
\textbf{偏差-方差权衡：}
\begin{itemize}[leftmargin=1.5em]
  \item $G_t$ 是 $v_\pi(S_t)$ 的无偏估计
  \item 真实 TD 目标 $R_{t+1}+\gamma v_\pi(S_{t+1})$ 也是无偏估计
  \item 实际使用的 TD 目标 $R_{t+1}+\gamma V(S_{t+1})$ 是有偏估计
  \item TD 目标的方差通常比 MC 的回报 $G_t$ 小很多
\end{itemize}
\end{keybox}

\notespace[2cm]

\subsection{时间差分学习的优势}

\begin{itemize}[leftmargin=2em]
  \item TD 不需要环境模型
  \item 与 MC 相比，TD 只需要等待一个时间步
  \item TD 只评估当前动作，与后继动作序列无直接关系
  \item TD 方法的收敛性已经得到证明
\end{itemize}

\notespace[2cm]

\subsection{表格型强化学习算法}

\subsubsection*{Sarsa：On-Policy TD}

学习行为值函数 $Q(s,a)$，基本数据单元是五元组 $(s,a,r,s',a')$。行动策略和评估策略都是同一个 $\varepsilon$-贪婪策略。

\textbf{更新公式：}

\begin{equation}
Q(s,a)\leftarrow Q(s,a)+\alpha\bigl[r+\gamma Q(s',a')-Q(s,a)\bigr]
\end{equation}

\subsubsection*{Q-learning：Off-policy TD}

直接逼近最优行为值函数。行动策略通常为 $\varepsilon$-贪婪，目标策略为贪婪。

\textbf{更新公式：}

\begin{equation}
Q(s,a)\leftarrow Q(s,a)+\alpha\bigl[r+\gamma \max_{a'}Q(s',a')-Q(s,a)\bigr]
\end{equation}

\subsubsection*{Expected Sarsa}

把下一状态下动作值的随机采样，替换成对策略分布下的期望，消除部分方差。

\subsubsection*{Double Q-learning}

\textbf{问题：} Q-learning 中的 $\max_{a'}Q(s',a')$ 会因最优化操作导致最大化偏差。

\textbf{思想：} 使用两个值函数逼近器，一个负责选择动作，另一个负责评估该动作的值，减轻过高估计问题。

\notespace[2cm]

\subsection{$\varepsilon$-贪婪策略}

\begin{equation}
\pi(a|s)= \begin{cases} 1-\varepsilon+\dfrac{\varepsilon}{|A(s)|}, & a=\argmax_a Q(s,a) \\ \dfrac{\varepsilon}{|A(s)|}, & \text{其他动作} \end{cases}
\end{equation}

\notespace[3cm]

\newpage

%==================================================================
% 第八章 函数逼近方法
%==================================================================
\section{函数逼近方法}

\subsection{表格型方法的局限}

当状态空间很大时，表格型方法存在以下问题：

\begin{itemize}[leftmargin=2em]
  \item 存储空间过大
  \item 学习速度慢
  \item 无法处理连续状态空间
\end{itemize}

\begin{keybox}
\textbf{解决方案：} 使用参数化函数逼近值函数或策略。
\begin{equation}
\hat{v}(s,\mathbf{w})\approx v_\pi(s)
\end{equation}
其中 $\mathbf{w}$ 是参数向量。
\end{keybox}

\notespace[2cm]

\subsection{线性函数逼近}

\subsubsection*{特征向量表示}

状态 $s$ 的特征向量：
\begin{equation}
\mathbf{x}(s)=(x_1(s),x_2(s),\dots,x_d(s))^T
\end{equation}

值函数的线性逼近：
\begin{equation}
\hat{v}(s,\mathbf{w})=\mathbf{w}^T\mathbf{x}(s)=\sum_{i=1}^{d}w_i x_i(s)
\end{equation}

\subsubsection*{目标函数}

最小化均方误差：
\begin{equation}
\mathrm{MSE}(\mathbf{w})=\sum_{s\in S}\bigl[v_\pi(s)-\hat{v}(s,\mathbf{w})\bigr]^2
\end{equation}

\subsubsection*{梯度下降更新}

\begin{equation}
\mathbf{w}_{t+1}=\mathbf{w}_t-\alpha\bigl[\hat{v}(s_t,\mathbf{w}_t)-v_\pi(s_t)\bigr]\nabla_{\mathbf{w}}\hat{v}(s_t,\mathbf{w}_t)
\end{equation}

对于线性逼近：
\begin{equation}
\nabla_{\mathbf{w}}\hat{v}(s,\mathbf{w})=\mathbf{x}(s)
\end{equation}

\begin{equation}
\mathbf{w}_{t+1}=\mathbf{w}_t+\alpha\bigl[v_\pi(s_t)-\hat{v}(s_t,\mathbf{w}_t)\bigr]\mathbf{x}(s_t)
\end{equation}

\notespace[2cm]

\subsection{蒙特卡洛与 TD 学习的函数逼近}

\subsubsection*{蒙特卡洛函数逼近}

使用采样回报 $G_t$ 作为目标：
\begin{equation}
\mathbf{w}_{t+1}=\mathbf{w}_t+\alpha\bigl[G_t-\hat{v}(s_t,\mathbf{w}_t)\bigr]\nabla_{\mathbf{w}}\hat{v}(s_t,\mathbf{w}_t)
\end{equation}

\subsubsection*{TD(0) 函数逼近}

使用 TD 目标：
\begin{equation}
\mathbf{w}_{t+1}=\mathbf{w}_t+\alpha\bigl[R_{t+1}+\gamma\hat{v}(s_{t+1},\mathbf{w}_t)-\hat{v}(s_t,\mathbf{w}_t)\bigr]\nabla_{\mathbf{w}}\hat{v}(s_t,\mathbf{w}_t)
\end{equation}

\begin{notebox}
\textbf{半梯度方法：} TD 方法中目标值也依赖于参数 $\mathbf{w}$，因此不是真正的梯度下降，称为半梯度方法。
\end{notebox}

\notespace[2cm]

\subsection{批量方法}

\subsubsection*{最小二乘蒙特卡洛}

给定一组样本 $\{(s_i,G_i)\}_{i=1}^{N}$，求解：
\begin{equation}
\mathbf{w}^*=\argmin_{\mathbf{w}}\sum_{i=1}^{N}\bigl[G_i-\mathbf{w}^T\mathbf{x}(s_i)\bigr]^2
\end{equation}

解析解：
\begin{equation}
\mathbf{w}^*=(\mathbf{X}^T\mathbf{X})^{-1}\mathbf{X}^T\mathbf{G}
\end{equation}

\subsubsection*{线性 TD 的收敛性}

TD(0) 的线性逼近收敛到：
\begin{equation}
\mathbf{w}_{TD}=\mathbf{A}^{-1}\mathbf{b}
\end{equation}
其中：
\begin{equation}
\mathbf{A}=\E\bigl[\mathbf{x}(s_t)(\mathbf{x}(s_t)-\gamma\mathbf{x}(s_{t+1}))^T\bigr]
\end{equation}
\begin{equation}
\mathbf{b}=\E\bigl[R_{t+1}\mathbf{x}(s_t)\bigr]
\end{equation}

\notespace[2cm]

\subsection{特征构造}

\subsubsection*{多项式基}

\begin{equation}
x_i(s)=s^i
\end{equation}

\subsubsection*{傅里叶基}

\begin{equation}
x_i(s)=\cos(i\pi s)
\end{equation}

\subsubsection*{径向基函数（RBF）}

\begin{equation}
x_i(s)=\exp\left(-\frac{\|s-c_i\|^2}{2\sigma_i^2}\right)
\end{equation}

\notespace[3cm]

\newpage

%==================================================================
% 第九章 神经网络与深度学习框架
%==================================================================
\section{神经网络与深度学习框架}

\subsection{神经网络基础}

\subsubsection*{神经元模型}

单个神经元的输出：
\begin{equation}
y=f\left(\sum_{i=1}^{n}w_i x_i+b\right)
\end{equation}
其中 $f$ 为激活函数，$w_i$ 为权重，$b$ 为偏置。

\subsubsection*{常见激活函数}

\begin{itemize}[leftmargin=2em]
  \item \textbf{Sigmoid}：$f(x)=\dfrac{1}{1+e^{-x}}$
  \item \textbf{ReLU}：$f(x)=\max(0,x)$
  \item \textbf{Tanh}：$f(x)=\dfrac{e^x-e^{-x}}{e^x+e^{-x}}$
\end{itemize}

\notespace[2cm]

\subsection{多层神经网络}

\subsubsection*{前向传播}

第 $l$ 层的输出：
\begin{equation}
\mathbf{a}^{(l)}=f\bigl(\mathbf{W}^{(l)}\mathbf{a}^{(l-1)}+\mathbf{b}^{(l)}\bigr)
\end{equation}

\subsubsection*{反向传播}

通过链式法则计算梯度：
\begin{equation}
\frac{\partial L}{\partial \mathbf{W}^{(l)}}=\frac{\partial L}{\partial \mathbf{a}^{(l)}}\cdot\frac{\partial \mathbf{a}^{(l)}}{\partial \mathbf{W}^{(l)}}
\end{equation}

\notespace[2cm]

\subsection{卷积神经网络（CNN）}

\subsubsection*{卷积操作}

\begin{equation}
(\mathbf{I}*\mathbf{K})_{i,j}=\sum_{m}\sum_{n}\mathbf{I}_{i+m,j+n}\mathbf{K}_{m,n}
\end{equation}

\subsubsection*{池化操作}

\begin{itemize}[leftmargin=2em]
  \item 最大池化：取窗口内最大值
  \item 平均池化：取窗口内平均值
\end{itemize}

\subsubsection*{CNN 的优势}

\begin{keybox}
\begin{itemize}[leftmargin=1.5em]
  \item 局部连接：减少参数数量
  \item 权值共享：进一步减少参数
  \item 平移不变性：对位置变化鲁棒
\end{itemize}
\end{keybox}

\notespace[2cm]

\subsection{深度学习框架}

\subsubsection*{TensorFlow}

\begin{itemize}[leftmargin=2em]
  \item Google 开发的开源框架
  \item 支持静态图和动态图
  \item 丰富的 API 和工具
\end{itemize}

\subsubsection*{PyTorch}

\begin{itemize}[leftmargin=2em]
  \item Facebook 开发的开源框架
  \item 动态图设计，更灵活
  \item 与 NumPy 接口相似，易于上手
\end{itemize}

\subsubsection*{基本操作示例}

\textbf{PyTorch 示例：}

\begin{verbatim}
import torch
x = torch.tensor([1.0, 2.0, 3.0])
y = torch.randn(3, 4)
z = torch.matmul(x, y)
\end{verbatim}

\notespace[2cm]

\subsection{优化算法}

\subsubsection*{随机梯度下降（SGD）}

\begin{equation}
\mathbf{w}_{t+1}=\mathbf{w}_t-\alpha\nabla_{\mathbf{w}}L(\mathbf{w}_t)
\end{equation}

\subsubsection*{Adam}

结合动量和自适应学习率：
\begin{equation}
\mathbf{m}_t=\beta_1\mathbf{m}_{t-1}+(1-\beta_1)\nabla L
\end{equation}
\begin{equation}
\mathbf{v}_t=\beta_2\mathbf{v}_{t-1}+(1-\beta_2)(\nabla L)^2
\end{equation}
\begin{equation}
\mathbf{w}_{t+1}=\mathbf{w}_t-\alpha\frac{\mathbf{m}_t}{\sqrt{\mathbf{v}_t}+\epsilon}
\end{equation}

\notespace[3cm]

\newpage

%==================================================================
% 第十章 DQN 及其变种
%==================================================================
\section{DQN 及其变种}

\subsection{DQN 的提出}

\begin{keybox}
\textbf{核心思想：} 使用深度神经网络逼近 Q 函数，结合经验回放和目标网络，实现稳定学习。
\end{keybox}

\subsubsection*{面临的挑战}

\begin{itemize}[leftmargin=2em]
  \item 数据相关性：强化学习顺序采集的数据高度相关
  \item 目标不稳定：目标值依赖于当前网络参数
  \item 过估计问题：最大化操作导致 Q 值过高估计
\end{itemize}

\notespace[2cm]

\subsection{经验回放}

\textbf{基本思想：} 将每次交互的经验 $(s,a,r,s')$ 存入经验池，训练时随机采样。

\begin{keybox}
\textbf{经验回放的作用：}
\begin{itemize}[leftmargin=1.5em]
  \item 打破数据相关性，使数据更接近独立同分布
  \item 提高数据利用效率，每个经验可多次使用
  \item 支持批量训练，提高训练稳定性
\end{itemize}
\end{keybox}

\notespace[2cm]

\subsection{目标网络}

\textbf{基本思想：} 使用独立的目标网络计算 TD 目标，定期更新目标网络参数。

\textbf{损失函数：}
\begin{equation}
L(\theta)=\E\bigl[\bigl(r+\gamma\max_{a'}Q(s',a';\theta^-)-Q(s,a;\theta)\bigr)^2\bigr]
\end{equation}
其中 $\theta^-$ 为目标网络参数。

\textbf{更新方式：}
\begin{itemize}[leftmargin=2em]
  \item 定期更新：每隔 $C$ 步将 $\theta^- \leftarrow \theta$
  \item 软更新：$\theta^- \leftarrow \tau\theta+(1-\tau)\theta^-$
\end{itemize}

\notespace[2cm]

\subsection{Double DQN}

\textbf{问题：} 标准 DQN 存在过估计问题。

\textbf{解决方案：} 使用两个网络，一个负责选择动作，另一个负责评估值。

\textbf{目标值计算：}
\begin{equation}
y=r+\gamma Q\bigl(s',\argmax_{a'}Q(s',a';\theta);\theta^- \bigr)
\end{equation}

\begin{notebox}
\textbf{区别：} 标准 DQN 用同一网络选择和评估；Double DQN 分离选择和评估，减少过估计。
\end{notebox}

\notespace[2cm]

\subsection{优先经验回放}

\textbf{基本思想：} 根据 TD 误差大小优先采样重要经验。

\textbf{优先级定义：}
\begin{equation}
p_i=|\delta_i|+\epsilon
\end{equation}

\textbf{采样概率：}
\begin{equation}
P(i)=\frac{p_i^\alpha}{\sum_k p_k^\alpha}
\end{equation}

\textbf{重要性权重：}
\begin{equation}
w_i=\left(\frac{1}{N\cdot P(i)}\right)^\beta
\end{equation}

\notespace[2cm]

\subsection{Dueling DQN}

\textbf{基本思想：} 将 Q 函数分解为状态值函数和优势函数。

\textbf{网络结构：}
\begin{equation}
Q(s,a)=V(s)+A(s,a)-\frac{1}{|A|}\sum_{a'}A(s,a')
\end{equation}

其中：
\begin{itemize}[leftmargin=2em]
  \item $V(s)$：状态值函数流
  \item $A(s,a)$：优势函数流
\end{itemize}

\begin{keybox}
\textbf{优势：} 对于状态值相同但动作值不同的场景，学习效率更高；对动作无关的状态变化更鲁棒。
\end{keybox}

\notespace[3cm]

\newpage

%==================================================================
% 第十一章 策略搜索方法
%==================================================================
\section{策略搜索方法}

\subsection{基于值函数的方法的局限}

\begin{keybox}
\textbf{局限性：}
\begin{itemize}[leftmargin=1.5em]
  \item 需要先学习值函数，再从值函数推导策略
  \item 当动作空间很大或连续时，$\argmax_a Q(s,a)$ 计算困难
  \item 策略可能过早收敛到次优解
\end{itemize}
\end{keybox}

\subsection{直接策略搜索的优势}

\begin{itemize}[leftmargin=2em]
  \item 直接在策略空间中优化，不需要中间的值函数
  \item 自然适用于连续动作空间
  \item 可以学习随机策略，有利于探索
\end{itemize}

\notespace[2cm]

\subsection{策略参数化}

策略由参数 $\theta$ 决定：$\pi_\theta(a|s)$。

\subsubsection*{离散动作空间}

常用玻尔兹曼分布：
\begin{equation}
\pi_\theta(a|s)=\frac{\exp\bigl(h(s,a;\theta)\bigr)}{\sum_{b}\exp\bigl(h(s,b;\theta)\bigr)}
\end{equation}

\subsubsection*{连续动作空间}

常用高斯策略：
\begin{equation}
\pi_\theta(a|s)=\frac{1}{\sqrt{2\pi\sigma^2}}\exp\left(-\frac{(a-\mu(s;\theta))^2}{2\sigma^2}\right)
\end{equation}

\notespace[2cm]

\subsection{目标函数}

\subsubsection*{目标函数定义}

\begin{equation}
J(\theta)=\E_{\pi_\theta}\bigl[v_{\pi_\theta}(s_0)\bigr]
\end{equation}

\subsubsection*{目标函数的梯度}

\begin{keybox}
\textbf{策略梯度定理：}
\begin{equation}
\nabla_\theta J(\theta)=\E_{\pi_\theta}\bigl[\nabla_\theta\log\pi_\theta(a|s)\,Q^{\pi_\theta}(s,a)\bigr]
\end{equation}
\end{keybox}

\notespace[2cm]

\subsection{REINFORCE 算法}

\textbf{基本思想：} 用采样回报 $G_t$ 替代 $Q^{\pi_\theta}(s,a)$。

\textbf{更新公式：}
\begin{equation}
\theta_{t+1}=\theta_t+\alpha\,G_t\,\nabla_\theta\log\pi_\theta(a_t|s_t)
\end{equation}

\begin{notebox}
\textbf{特点：} 无偏估计，但方差较大；属于蒙特卡洛方法，需要完整 episode。
\end{notebox}

\notespace[2cm]

\subsection{Actor-Critic 方法}

\textbf{基本思想：} 使用学习到的值函数作为 Critic，减少方差。

\textbf{更新公式：}
\begin{equation}
\theta_{t+1}=\theta_t+\alpha\,Q_w(s_t,a_t)\,\nabla_\theta\log\pi_\theta(a_t|s_t)
\end{equation}

其中 $Q_w$ 为 Critic 网络逼近的行为值函数。

\begin{keybox}
\textbf{优势：}
\begin{itemize}[leftmargin=1.5em]
  \item 方差比 REINFORCE 小
  \item 可以在线学习，不需要完整 episode
  \item 结合了值函数方法和策略梯度方法的优势
\end{itemize}
\end{keybox}

\notespace[3cm]

\newpage

%==================================================================
% 第十二章 随机策略梯度
%==================================================================
\section{随机策略梯度}

\subsection{策略梯度定理}

\begin{keybox}
\textbf{核心定理：}
\begin{equation}
\nabla_\theta J(\theta)=\int_S\rho^{\pi}(s)\int_A\nabla_\theta\pi_\theta(a|s)Q^{\pi}(s,a)\,da\,ds
\end{equation}
\end{keybox}

\subsubsection*{等价形式}

利用 $\nabla_\theta\pi_\theta(a|s)=\pi_\theta(a|s)\nabla_\theta\log\pi_\theta(a|s)$：

\begin{equation}
\nabla_\theta J(\theta)=\E_{s\sim\rho^{\pi},a\sim\pi_\theta}\bigl[\nabla_\theta\log\pi_\theta(a|s)\,Q^{\pi}(s,a)\bigr]
\end{equation}

\notespace[2cm]

\subsection{蒙特卡洛策略梯度（REINFORCE）}

\textbf{算法流程：}

\begin{algorithmic}[1]
\For{每个 episode}
  \State 按 $\pi_\theta$ 生成轨迹 $s_0,a_0,r_1,\dots,s_{T-1},a_{T-1},r_T$
  \For{$t=0$ 到 $T-1$}
    \State $G_t=\sum_{k=t}^{T-1}\gamma^{k-t}r_{k+1}$
    \State $\theta\leftarrow\theta+\alpha\gamma^t G_t\nabla_\theta\log\pi_\theta(a_t|s_t)$
  \EndFor
\EndFor
\end{algorithmic}

\begin{notebox}
\textbf{特点：} 使用实际回报 $G_t$ 作为 $Q^{\pi}(s,a)$ 的估计；无偏但方差大。
\end{notebox}

\notespace[2cm]

\subsection{Actor-Critic 结构}

\subsubsection*{基本框架}

\begin{itemize}[leftmargin=2em]
  \item \textbf{Actor}：策略网络 $\pi_\theta(a|s)$，负责选择动作
  \item \textbf{Critic}：值函数网络 $V_w(s)$ 或 $Q_w(s,a)$，负责评估动作
\end{itemize}

\subsubsection*{Actor 更新}

\begin{equation}
\theta\leftarrow\theta+\alpha\nabla_\theta\log\pi_\theta(a|s)\,A(s,a)
\end{equation}

其中 $A(s,a)=Q(s,a)-V(s)$ 为优势函数。

\subsubsection*{Critic 更新}

TD 学习：
\begin{equation}
w\leftarrow w+\alpha\bigl[r+\gamma V_w(s')-V_w(s)\bigr]\nabla_w V_w(s)
\end{equation}

\notespace[2cm]

\subsection{优势函数估计}

\subsubsection*{优势函数的作用}

\begin{keybox}
\textbf{优势函数：} $A(s,a)=Q(s,a)-V(s)$

\textbf{含义：} 动作 $a$ 相对于平均水平的优劣程度。使用优势函数可以减少方差，提高学习稳定性。
\end{keybox}

\subsubsection*{优势函数估计方法}

\textbf{1. TD 误差：}
\begin{equation}
\delta_t=r_t+\gamma V(s_{t+1})-V(s_t)
\end{equation}

\textbf{2. GAE（Generalized Advantage Estimation）：}
\begin{equation}
A_t=\sum_{l=0}^{\infty}(\gamma\lambda)^l\delta_{t+l}
\end{equation}

其中 $\lambda$ 为平衡参数。

\notespace[2cm]

\subsection{相容函数}

\begin{keybox}
\textbf{相容条件：} 若 Critic 满足以下条件，则策略梯度无偏：
\begin{enumerate}[leftmargin=1.5em]
  \item $Q_w(s,a)=\nabla_\theta\log\pi_\theta(a|s)^T w$
  \item 参数 $w$ 使均方误差最小
\end{enumerate}
\end{keybox}

\textbf{相容函数形式：}
\begin{equation}
Q_w(s,a)=\nabla_\theta\log\pi_\theta(a|s)^T w
\end{equation}

\notespace[2cm]

\subsection{随机 Off-policy Actor-Critic}

\textbf{采样策略：} $\beta(a|s)$

\textbf{目标策略：} $\pi_\theta(a|s)$

\textbf{目标函数：}
\begin{equation}
J_\beta(\pi_\theta)=\int_S\rho^\beta(s)V^\pi(s)\,ds
\end{equation}

\textbf{重要性采样修正：}
\begin{equation}
\nabla_\theta J(\theta)=\E_{s\sim\rho^\beta,a\sim\beta}\bigl[\frac{\pi_\theta(a|s)}{\beta(a|s)}\nabla_\theta\log\pi_\theta(a|s)Q^\pi(s,a)\bigr]
\end{equation}

\notespace[3cm]

\newpage

%==================================================================
% 第十三章 TRPO 与 PPO
%==================================================================
\section{TRPO 与 PPO}

\subsection{策略梯度算法的缺点}

\begin{keybox}
\textbf{主要问题：}
\begin{itemize}[leftmargin=1.5em]
  \item 步长选择困难
  \item 策略过早收敛到次优确定性策略
  \item 步长不合适时学习过程可能崩溃
\end{itemize}
\end{keybox}

\notespace[2cm]

\subsection{替代回报函数}

\textbf{基本思想：} 寻找一个易于优化的替代目标函数。

\subsubsection*{优势函数}

\begin{equation}
A^{\pi_{\text{old}}}(s,a)=Q^{\pi_{\text{old}}}(s,a)-V^{\pi_{\text{old}}}(s)
\end{equation}

\subsubsection*{替代目标函数}

\begin{equation}
L_{\pi_{\text{old}}}(\pi)=\eta(\pi_{\text{old}})+\sum_s\rho_{\pi_{\text{old}}}(s)\sum_a\pi(a|s)A^{\pi_{\text{old}}}(s,a)
\end{equation}

\subsubsection*{重要性采样形式}

\begin{equation}
L_{\theta_{\text{old}}}(\theta)=\hat{\E}_t\left[\frac{\pi_\theta(a_t|s_t)}{\pi_{\theta_{\text{old}}}(a_t|s_t)}\hat{A}_t\right]
\end{equation}

\notespace[2cm]

\subsection{理论下界}

\begin{keybox}
\textbf{定理：} 令 $\alpha=D_{\max}^{TV}(\pi_{\text{old}},\pi_{\text{new}})$，则
\begin{equation}
\eta(\pi_{\text{new}})\ge L_{\pi_{\text{old}}}(\pi_{\text{new}})-\frac{4\epsilon\gamma}{(1-\gamma)^2}\alpha^2
\end{equation}
其中 $\epsilon=\max_{s,a}|A^{\pi_{\text{old}}}(s,a)|$。
\end{keybox}

\textbf{含义：} 若新旧策略差异过大，替代目标对真实回报的近似会变差。

\notespace[2cm]

\subsection{TRPO 实用算法}

\textbf{优化问题：}
\begin{equation}
\max_\theta L_{\theta_{\text{old}}}(\theta) \quad \text{s.t.} \quad \bar{D}_{KL}^{\rho_{\theta_{\text{old}}}(\theta_{\text{old}},\theta)}\le\delta
\end{equation}

\subsubsection*{关键技巧}

\begin{enumerate}[leftmargin=2em]
  \item 忽略状态分布变化
  \item 使用重要性采样
  \item 利用平均 KL 散度替代逐状态约束
  \item 利用当前策略对应的状态分布
\end{enumerate}

\subsubsection*{共轭梯度法}

求解自然梯度方向：$x=A^{-1}g$，其中 $A$ 为 Fisher 信息矩阵。

\notespace[2cm]

\subsection{PPO 算法}

\subsubsection*{CLIP 目标函数}

\begin{keybox}
\begin{equation}
L^{CLIP}(\theta)=\hat{\E}_t\left[\min\left(r_t(\theta)\hat{A}_t,\,\text{clip}\bigl(r_t(\theta),1-\epsilon,1+\epsilon\bigr)\hat{A}_t\right)\right]
\end{equation}
其中 $r_t(\theta)=\dfrac{\pi_\theta(a_t|s_t)}{\pi_{\theta_{\text{old}}}(a_t|s_t)}$。
\end{keybox}

\textbf{含义：}
\begin{itemize}[leftmargin=2em]
  \item 当 $\hat{A}_t>0$ 时，若 $r_t$ 超过 $1+\epsilon$，则限制其继续增大
  \item 当 $\hat{A}_t<0$ 时，若 $r_t$ 小于 $1-\epsilon$，则限制其继续减小
\end{itemize}

\subsubsection*{PPO 总目标函数}

\begin{equation}
L(\theta)=\hat{\E}_t\left[L_t^{CLIP}(\theta)-c_1 L_t^{VF}(\theta)+c_2 S[\pi_\theta](s_t)\right]
\end{equation}

其中 $L^{VF}$ 为值函数损失，$S$ 为策略熵。

\notespace[2cm]

\subsection{回报函数对比}

\begin{summarybox}
\textbf{策略梯度：}
\begin{equation}
L^{PG}(\theta)=\hat{\E}_t\left[\log\pi_\theta(a_t|s_t)\hat{A}_t\right]
\end{equation}

\textbf{TRPO：}
\begin{equation}
L^{TRPO}(\theta)=\hat{\E}_t\left[\frac{\pi_\theta(a_t|s_t)}{\pi_{\theta_{\text{old}}}(a_t|s_t)}\hat{A}_t\right]
\end{equation}
并加入 KL 约束。

\textbf{PPO：}
\begin{equation}
L^{CLIP}(\theta)=\hat{\E}_t\left[\min\left(r_t(\theta)\hat{A}_t,\,\text{clip}\bigl(r_t(\theta),1-\epsilon,1+\epsilon\bigr)\hat{A}_t\right)\right]
\end{equation}
\end{summarybox}

\notespace[2cm]

\subsection{信息论基础}

\subsubsection*{熵}

离散系统熵：
\begin{equation}
H(X)=-\sum_x p(x)\log p(x)
\end{equation}

\subsubsection*{KL 散度}

\begin{equation}
D_{KL}(p||q)=\sum_x p(x)\log\frac{p(x)}{q(x)}
\end{equation}

\begin{notebox}
\textbf{在 TRPO 中：} KL 散度用于约束新旧策略之间的差异，控制更新步长。
\end{notebox}

\notespace[3cm]

\newpage

%==================================================================
% 第十四章 确定性策略梯度
%==================================================================
\section{确定性策略梯度}

\subsection{随机策略与确定性策略}

\begin{keybox}
\textbf{区别：}
\begin{itemize}[leftmargin=1.5em]
  \item \textbf{随机策略}：$a\sim\pi_\theta(a|s)$，需要对状态空间和动作空间积分
  \item \textbf{确定性策略}：$a=\mu_\theta(s)$，只需对状态空间积分
\end{itemize}
\textbf{当动作空间维数较高时，确定性策略方法更有优势。}
\end{keybox}

\subsubsection*{随机策略的参数化}

随机策略可写为确定性策略 + 随机扰动：
\begin{equation}
a=\mu_\theta(s)+\epsilon, \quad \epsilon\sim\mathcal{N}(0,\sigma^2)
\end{equation}

即 $a\sim\mathcal{N}(\mu_\theta(s),\sigma^2)$。

\notespace[2cm]

\subsection{确定性值函数}

行为值函数：
\begin{equation}
Q^\pi(s_t,a_t)=\E\left[\sum_{i=t}^{\infty}\gamma^{i-t}r_i\,\middle|\, s_t,a_t\right]
\end{equation}

迭代形式：
\begin{equation}
Q^\pi(s_t,a_t)=\E\left[r(s_t,a_t)+\gamma Q^\pi(s_{t+1},a_{t+1})\right]
\end{equation}

若目标策略为确定性策略 $\mu$：
\begin{equation}
Q^\mu(s_t,a_t)=\E\left[r(s_t,a_t)+\gamma Q^\mu(s_{t+1},\mu(s_{t+1}))\right]
\end{equation}

\notespace[2cm]

\subsection{确定性策略梯度定理}

\begin{keybox}
\textbf{核心思想：} 根据链式法则，梯度通过 $Q^\mu(s,a)$ 对动作的导数，以及动作对参数 $\theta$ 的导数传递。

\textbf{策略梯度定理：}
\begin{equation}
\nabla_\theta J(\mu_\theta)=\E_{s\sim\rho^\mu}\bigl[\nabla_\theta\mu_\theta(s)\nabla_a Q^\mu(s,a)\big|_{a=\mu_\theta(s)}\bigr]
\end{equation}
\end{keybox}

\notespace[2cm]

\subsection{相容函数逼近}

\subsubsection*{条件 1：动作梯度匹配}

\begin{equation}
\nabla_a Q_w(s,a)\big|_{a=\mu_\theta(s)}=\nabla_\theta\mu_\theta(s)^T w
\end{equation}

\subsubsection*{条件 2：最小化均方误差}

\begin{equation}
\mathrm{MSE}(w)=\E\bigl[\varepsilon(s;\theta,w)^T\varepsilon(s;\theta,w)\bigr]
\end{equation}

其中 $\varepsilon(s;\theta,w)=\nabla_a Q_w(s,a)\big|_{a=\mu_\theta(s)}-\nabla_a Q^\mu(s,a)\big|_{a=\mu_\theta(s)}$。

\subsubsection*{满足条件 1 的相容函数形式}

\begin{equation}
Q_w(s,a)=(a-\mu_\theta(s))^T\nabla_\theta\mu_\theta(s)^T w+V^\mu(s)
\end{equation}

\notespace[2cm]

\subsection{DDPG 算法}

\subsubsection*{面临的问题}

\begin{itemize}[leftmargin=2em]
  \item 神经网络训练假设数据独立同分布，而强化学习顺序采集的数据存在相关性
  \item 若能采用 mini-batch 优化，可更充分利用数据
\end{itemize}

\subsubsection*{DQN 的成功经验}

经验回放（Experience Replay）、独立目标网络（Target Network）。

\begin{keybox}
\textbf{DDPG 的核心机制：}
\begin{itemize}[leftmargin=1.5em]
  \item 行动策略为随机策略
  \item 使用目标网络
  \item 使用经验回放
  \item 目标网络参数进行延迟/软更新
  \item 采用 Actor-Critic 结构处理连续动作问题
\end{itemize}
\end{keybox}

\notespace[2cm]

\subsection{TD3：DDPG 的扩展}

\subsubsection*{DDPG 存在的问题}

\begin{itemize}[leftmargin=2em]
  \item 值函数过高估计
  \item 值函数方差较高
  \item 策略网络与值函数网络之间存在耦合关系
\end{itemize}

\subsubsection*{解决思路}

\begin{itemize}[leftmargin=2em]
  \item 使用独立目标网络
  \item \textbf{Twin Delayed 机制}：目标网络更新延迟、策略更新延迟、目标策略平滑
\end{itemize}

\notespace[2cm]

\subsection{D4PG}

\begin{notebox}
\textbf{全称：} Distributed Distributional Deep Deterministic Policy Gradient

\textbf{核心思想：} 任何改善 Critic 学习质量的方法，都会直接改善 Actor 更新质量。值分布建模了函数逼近等带来的随机性，利用值分布更新 Critic 可以直接提升算法表现。
\end{notebox}

\subsubsection*{D4PG 使用的关键技巧}

\begin{itemize}[leftmargin=2em]
  \item 分布式值函数强化学习，得到更细致的 Critic
  \item 分布式并行 Actor
  \item $N$-step 回报
  \item 优先经验回放采样
\end{itemize}

\notespace[3cm]

\newpage

%==================================================================
% 第十五章 博弈论与多智能体强化学习
%==================================================================
\section{博弈论与多智能体强化学习}

\subsection{多智能体系统与分类}

\subsubsection*{按任务类型分类}

\begin{itemize}[leftmargin=2em]
  \item 完全协作多智能体系统
  \item 完全竞争多智能体系统
  \item 混合策略多智能体系统
\end{itemize}

\subsubsection*{多智能体强化学习的挑战}

\begin{keybox}
\begin{itemize}[leftmargin=1.5em]
  \item MDP 假设环境是静态的；但从单智能体视角看，其他智能体本身是动态环境，因此环境不再平稳
  \item 当智能体之间是竞争关系时，最优策略往往不存在，需要求解均衡解
\end{itemize}
\end{keybox}

\notespace[2cm]

\subsection{博弈论基本分类}

\subsubsection*{按参与人行动先后顺序}

\begin{itemize}[leftmargin=2em]
  \item \textbf{静态博弈}：参与人同时选择行动（如猜丁壳），常用标准型表述（矩阵博弈）
  \item \textbf{动态博弈}：参与人的动作有先后顺序，且后行动者可观察前行动者的选择（如围棋），常用扩展型表述（博弈树）
\end{itemize}

\subsubsection*{按参与人对信息掌握程度}

\begin{itemize}[leftmargin=2em]
  \item \textbf{完美信息博弈}：参与人对相关信息已知（如棋类游戏）
  \item \textbf{不完美信息博弈}：参与人对相关信息并不完全知道（如牌类游戏）
\end{itemize}

\notespace[2cm]

\subsection{完美信息静态博弈的基本概念}

\subsubsection*{基本要素}

\begin{itemize}[leftmargin=2em]
  \item \textbf{参与者}：参与博弈的智能体 $i$
  \item \textbf{动作空间}：$A_i$，联合动作 $a=(a_1,a_2,\dots,a_n)$
  \item \textbf{策略}：$\sigma_i$
  \begin{itemize}
    \item 纯策略：选择某一行为的概率为 $1$
    \item 混合策略：选择某一行为的概率小于 $1$
  \end{itemize}
  \item \textbf{联合策略}：$\sigma=(\sigma_1,\sigma_2,\dots,\sigma_n)$
  \item \textbf{回报函数}：$R_i(a)$
  \item \textbf{值函数}：$V_i(\sigma)$
\end{itemize}

\notespace[2cm]

\subsection{最优响应与纳什均衡}

\subsubsection*{最优响应}

当其他玩家策略固定时，玩家 $k$ 的最优响应：
\begin{equation}
R_k(\sigma_{-k})=\bigl\{\sigma_k'\in A_k\,\big|\,\mu_k(\sigma_k',\sigma_{-k})\ge\mu_k(\sigma_k,\sigma_{-k}),\ \forall\sigma_k\in A_k\bigr\}
\end{equation}

\subsubsection*{纳什均衡}

\begin{keybox}
当所有玩家都采取各自的最优响应时，称该联合策略为纳什均衡。
\begin{equation}
V_i(\sigma_i^*,\sigma_{-i}^*)\ge V_i(\sigma_i,\sigma_{-i}^*), \quad \forall i,\ \forall\sigma_i
\end{equation}
\end{keybox}

\textbf{例子：} 囚徒困境中，$(\text{坦白},\text{坦白})$ 是纳什均衡策略。

\notespace[2cm]

\subsection{双人零和博弈}

双人联合行为对应的回报构成矩阵，因此称为矩阵博弈。

\begin{keybox}
求解双人零和矩阵博弈中的纳什均衡，等价于求解相应的极小极大问题，进一步可转化为线性规划问题，使用单纯形法求解。
\end{keybox}

\notespace[2cm]

\subsection{非完美信息博弈：扩展式博弈}

七元组描述：$(H,Z,P,p,u,I,\sigma)$

\begin{itemize}[leftmargin=2em]
  \item $H$：历史节点集合
  \item $Z$：终止状态集合
  \item $P$：玩家集合
  \item $p$：非终止状态到玩家的映射
  \item $u$：终止状态到实数的映射（回报）
  \item $I$：信息集，其中所有节点对玩家 $i$ 来说不可区分
  \item $\sigma$：策略
\end{itemize}

\begin{notebox}
\textbf{perfect recall（完美回放）：} 指玩家对于先前已知的信息不会发生任何程度的遗忘。
\end{notebox}

\notespace[2cm]

\subsection{随机博弈（马尔科夫博弈）}

随机博弈（Stochastic game）核心元素：$n$（玩家个数）、联合状态转移概率、联合立即回报、联合动作空间、局部策略、联合策略。

\subsubsection*{值函数特征}

在马尔科夫博弈中，每个智能体的行为值函数依赖于联合动作。

\notespace[2cm]

\subsection{多智能体强化学习}

\subsubsection*{回报定义}

对于智能体 $i$：
\begin{equation}
V_i^\pi(s)=\E_\pi\left\{\sum_{t=0}^{\infty}\gamma^t r_i(t+1)\,\middle|\, s(0)=s\right\}
\end{equation}
其中 $\pi=(\pi_1,\pi_2,\dots,\pi_n)$ 为联合策略。

\begin{keybox}
\textbf{注意：} 一般情况下，不可能同时最大化每个智能体的目标函数，因为不同智能体的目标可能相互矛盾。最优响应与纳什均衡的概念可扩展到马尔科夫博弈中。
\end{keybox}

\notespace[2cm]

\subsection{多智能体 Q-learning}

\subsubsection*{更新形式}

\begin{equation}
Q(s,a)\leftarrow(1-\alpha)Q(s,a)+\alpha\bigl[R+\gamma V(s')\bigr]
\end{equation}

\subsubsection*{下一个状态值函数 $V(s')$ 的获取方式}

\textbf{1. 对手建模方法：} 统计其他智能体可能联合动作的频率 $F(s,a_{-i})$：
\begin{equation}
V_i(s)=\max_{a_i\in A_i}\sum_{a_{-i}\in A_{-i}}F(s,a_{-i})\,Q_i(s,a_i,a_{-i})
\end{equation}

\textbf{2. 假设对手按固定规则选动作}（如零和博弈中最小最大原则）：
\begin{equation}
V(s)=\min_\mu\max_\sigma\sum_{a\in A}\sigma(a)\mu(a')Q(s,a,a')
\end{equation}

\textbf{3. 假设下一个状态执行均衡策略}（Nash-Q、相关均衡、Stackelberg 均衡）：
\begin{equation}
V(s)=\text{Nash}_s\,Q(s,\cdot)
\end{equation}

\notespace[2cm]

\subsection{多智能体强化学习算法类型}

\subsubsection*{1. 完全协作强化学习算法}

示例：Cooperative and Distributed Reinforcement Learning of Drones for Field Coverage。

\subsubsection*{2. 完全竞争：Minimax Q-learning}

定义值函数时采用最小最大思想。若相关博弈模型已知，可利用线性规划方法求解。

\subsubsection*{3. 混合任务：Nash-Q 学习}

\begin{itemize}[leftmargin=2em]
  \item 利用二次规划更新行为值函数
  \item 通过求解均衡策略进行更新
  \item 收敛条件：每个状态处的博弈都有一个全局最优点或者鞍点（过于严格，实际中很难满足）
\end{itemize}

\subsubsection*{4. 混合任务：朋友或敌人 Q-learning}

其他玩家要么被视为"朋友"（共同最大化），要么被视为"敌人"（共同最小化）。$N$ 玩家一般和随机博弈可被看作一个动作集扩展的两玩家零和博弈。

\notespace[2cm]

\subsection{WoLF 策略爬山算法}

\begin{keybox}
\textbf{基本思想：} WoLF-PHC 算法仅需要每个玩家自己的动作信息，可显著约减空间复杂度。组合了两个方法：
\begin{itemize}[leftmargin=1.5em]
  \item \textbf{Win or Learn Fast（WoLF）}：做得比期望差时快速学习，做得好时谨慎学习
  \item \textbf{Policy Hill-Climbing（PHC）}：策略爬山算法，满足合理性要求
\end{itemize}
\end{keybox}

\textbf{特点：} 策略爬山算法仅能保证在静态环境下，单智能体收敛到最优策略。

\notespace[2cm]

\subsection{基于微分对策的方法}

内容要点：运动学方程、性能指标函数、定义双人零和微分对策、推导 Hamilton-Jacobi-Isaacs 方程。

\notespace[2cm]

\subsection{MADDPG}

\begin{notebox}
\textbf{MADDPG 算法框架：} 多智能体深度确定性策略梯度方法，结合了 DDPG 和多智能体学习。

\textbf{核心思想：} 每个智能体有自己的 Actor-Critic 结构，Critic 可以观察其他智能体的策略。
\end{notebox}

\notespace[3cm]

\newpage

%==================================================================
% 附录：参考材料
%==================================================================
\section*{附录：参考材料}

\subsection*{主要参考书籍}

\begin{itemize}[leftmargin=2em]
  \item Sutton, Barto: \textit{Reinforcement Learning: An Introduction}
  \item 郭宪、方勇纯：《深入浅出强化学习：原理入门》
  \item 郭宪、宋俊潇、方勇纯：《深入浅出强化学习：编程实战》
\end{itemize}

\subsection*{主要参考论文}

\begin{itemize}[leftmargin=2em]
  \item DQN: Mnih et al., 2015
  \item Double DQN: van Hasselt et al., 2016
  \item Prioritized Experience Replay: Schaul et al., 2016
  \item Dueling DQN: Wang et al., 2016
  \item DDPG: Lillicrap et al., 2016
  \item TD3: Fujimoto et al., 2018
  \item TRPO: Schulman et al., 2015
  \item PPO: Schulman et al., 2017
  \item MADDPG: Lowe et al., 2017
\end{itemize}

\notespace[5cm]

\end{document}